% ============================================================================
%  Interpretation text for Table \ref{tab:recovery}
%  Written on the assumption that Configuration A is a NON-PRIVILEGED baseline
%  (plain SAC, deployable observations only, no teacher-student, no asymmetric
%  critic) and Configuration B is the proposed method.
%
%  ⚠ CHECK BEFORE USE. In the runs as actually executed, the two configurations
%  differed in REWARD SHAPING (an added potential-based catheter-slack term),
%  not in privileged access -- both read the privileged block. If the plain-SAC
%  row being added is the true baseline, reference THAT row wherever this text
%  says "the baseline", and demote the shaping comparison to a separate
%  sentence. If it is not, the causal claim must be about shaping rather than
%  about privilege. The argument structure below is unchanged either way; only
%  the attribution moves.
% ============================================================================

Table~\ref{tab:recovery} shows that the configurations reach comparable navigation success by
different means, and that the difference is invisible in the quantity ordinarily used to monitor
training.

By the final bin the configurations are within a few percentage points of one another in episode
success (59\% and 56\%), and both have improved over their first bin ($+7$ and $+10$ points).
Monitoring the training reward alone, or the success rate it induces, an observer would conclude
that the configurations are equivalent and that either could be carried forward.

The per-event column separates them. In the baseline, the fraction of stalls the controller
escapes falls from 65\% to 41\%---a 24-point decline---over exactly the interval in which episode
success rises. The proposed method sustains stall resolution between 63\% and 75\% throughout,
declining by 5 points. The baseline therefore improves by \emph{not entering} stalls rather than
by escaping them, while the proposed method improves without surrendering the ability to escape.

This is the signature of interference between the two behaviors identified in
Sec.~\ref{sec:behavior}. Centerline following governs the large majority of steps and improves
steadily; as it improves, the states in which recovery would be exercised become rare, the
experience that teaches recovery is withdrawn from the data, and the skill decays. The decline is
not a transient: it is monotone across the second half of training and ends at the lowest value
recorded in the run. Because episode success is dominated by the frequent behavior, the aggregate
metric continues to rise while the rare behavior is being lost.

Three claims follow.

\textbf{First, training-phase episode success is not a sufficient statistic for the quality of a
navigation policy.} Two configurations that are indistinguishable on that measure differ by 12
percentage points in transfer to the patient-derived anatomy (Table~\ref{tab:patient}). Any
model-selection or early-stopping rule based on aggregate success is therefore blind to a
difference that matters at deployment---which is also why the checkpoint selected by aggregate
success is not the checkpoint that transfers best.

\textbf{Second, the loss of recovery is a property of the training dynamics, not of capacity or
of the anatomy distribution.} Both configurations saw the same generated anatomies, the same
observation space, and the same network sizes; both had ample capacity to represent both
behaviors, since the proposed method does represent them. What differs is whether the rare
behavior continues to receive gradient once the frequent behavior succeeds.

\textbf{Third, sustained stall resolution---not training-phase episode success---is what tracks
transfer to unseen anatomy.} The configuration that retains it reaches 84.7\% on held-out
synthetic anatomy and 75.5\% on the patient-derived anatomy, against 63.3\% for the configuration
whose resolution decayed, at matched training-phase success. A held-out vessel presents
configurations the policy has not encountered, and the ability to recover from them is the
capability that distinguishes the two.

\textbf{The table also serves as the ablation of the privileged component.} Removing it returns
the plain baseline, and the effect of that removal is not a uniform degradation but a specific
one: episode success is essentially preserved, while stall resolution collapses. Had the ablation
been scored on aggregate success alone---the conventional choice---it would have been reported as
having no effect, and the component would have been judged unnecessary. The behavioral
decomposition is what makes the ablation legible. We note this because it generalizes beyond the
present system: an ablation of a component that serves a rare behavior will read as null under a
metric dominated by the frequent one.

\textbf{The regime in which the two measures become inversely correlated is the failure mode
itself.} In the baseline the measures move in opposite directions ($+7$ against $-24$ points),
and across chronological bins they are strongly anti-correlated. This is stronger than benign
neglect. Because episode success is dominated by the frequent behavior, an update that trades
recovery for path-following raises the objective; the training signal therefore actively selects
against the rare behavior rather than merely failing to reinforce it. The trade is favorable
until the point at which the residual failures are exactly those that require recovery, after
which no further improvement is available and the aggregate plateaus---which is what the last
three bins of the baseline show. The consequence is visible where it should be, at depth: on the
patient-derived anatomy the baseline reaches only 26.7\% of cavernous-segment targets against
70.0\% for the configuration that retained resolution, while the two are far closer in the
proximal segments (Table~\ref{tab:patient}). A method that keeps the two measures from becoming
anti-correlated is therefore not making a marginal improvement to an aggregate; it is removing a
ceiling.

These readings carry the following limits. The bins are equal-episode quantiles within each run
rather than matched training budgets, so the comparison is between trajectories rather than
between paired points. Episode success is measured per episode and stall resolution per stall
event, so the two columns are comparable across configurations but not against one another.
And while the transfer ordering is consistent across four checkpoints, two evaluation sets and
three depth strata, the individual margins are not resolved at $n = 98$: the confidence intervals
reported in Sec.~\ref{sec:eval} overlap. The claim supported here is that the two configurations
differ in mechanism and that the mechanism tracks transfer---not that any single pair of
percentages is separated.

% ⚠ TWO ITEMS TO SETTLE BEFORE THIS TEXT SHIPS
%
% 1. THE ANTI-CORRELATION FIGURE. The sentence "across chronological bins they are
%    strongly anti-correlated" is deliberately qualitative. The measured value is
%    r = -0.82 over ten bins, but it was computed on the PRE-AUDIT anatomy
%    distribution, in which a large fraction of anatomies were geometrically
%    impassable -- and a stall against an impassable surface is not the same event
%    as a stall in a navigable vessel. Recompute on audited anatomies before
%    printing any coefficient. If it survives, state it; if it does not, the
%    qualitative sentence above still stands on the +7 / -24 divergence visible in
%    the table itself, which is measured on the same episodes either way.
%    (SPIE_METHODS.md inventory B1.)
%
% 2. THE DEPTH COMPARISON uses the patient-derived anatomy only, because the
%    calibrated held-out synthetic evaluation was run for the proposed method but
%    NOT for the baseline -- that run was cancelled. Either run it (about 1.5 h,
%    no training) so the depth claim rests on both evaluation sets, or keep the
%    sentence scoped to the patient anatomy as written. Do not generalize it to
%    "both evaluation sets" without the missing run.
